% =============================================================================
%  Seccion: Datos
%  Contenido de EJEMPLO. Muestra listas y tablas con booktabs.
% =============================================================================

\section{Datos}\label{sec:datos}

De donde salieron los datos, cuantas observaciones y variables hay, y que
limpieza se les hizo. Si el dataset vive en \texttt{data/}, dilo con su nombre
exacto de archivo.

Las listas van asi:

\begin{itemize}
  \item \textbf{Origen:} archivo \texttt{data/ejemplo.csv}.
  \item \textbf{Unidad de observacion:} un renglon por individuo.
  \item \textbf{Periodo:} enero a diciembre de 2025.
\end{itemize}

Y cuando el orden importa, con \verb|enumerate|:

\begin{enumerate}
  \item Se eliminaron los registros duplicados.
  \item Se imputaron los faltantes con la mediana de cada columna.
  \item Se estandarizaron las variables a media cero y varianza uno.
\end{enumerate}

Las tablas se arman con \texttt{booktabs}: solo lineas horizontales
(\verb|\toprule|, \verb|\midrule|, \verb|\bottomrule|) y ninguna vertical. La
leyenda de una tabla va \emph{arriba}. Los resultados estan en
\cref{tab:descriptivos}.

\begin{table}[t]
  \centering
  \caption{Estadisticos descriptivos de las variables estandarizadas.}
  \label{tab:descriptivos}
  \begin{tabular}{lrrr}
    \toprule
    Variable & Media & Desv. est. & $n$ \\
    \midrule
    $X_1$ & 0.00 & 1.00 & 120 \\
    $X_2$ & 0.00 & 1.00 & 120 \\
    $X_3$ & 0.00 & 1.00 & 118 \\
    \bottomrule
  \end{tabular}
\end{table}

El \texttt{[t]} de \verb|\begin{table}[t]| es la posicion sugerida: \emph{top}
de la pagina. LaTeX mueve los flotantes a donde quepan, y eso es correcto: no
pelees con el poniendo \texttt{[H]} en todas partes, porque de ahi salen las
paginas medio vacias.
